
\section{Theoretic Results}
\label{sec:expressiveness}

\paragraph{Setup.}
Let a motion sample be represented by segment features
$\{f_{m_i}\}_{i=1}^m \subset \mathbb{R}^d$
and a text sample by token features
$\{f_{t_j}\}_{j=1}^n \subset \mathbb{R}^d$.
Define the pairwise similarity matrix
$C\in\mathbb{R}^{m\times n}$ by
\[
C_{ij}
\;=\;
\langle f_{m_i}, f_{t_j}\rangle .
\]
Let the uniform marginals be
$\mu = \frac{1}{m}\mathbf{1}_m$
and
$\nu = \frac{1}{n}\mathbf{1}_n$,
and define the transport polytope
\[
\Pi(\mu,\nu)
=
\left\{
\pi\in\mathbb{R}_+^{m\times n}
\;\middle|\;
\begin{aligned}
\pi\mathbf{1}_n &= \mu, \\
\pi^\top \mathbf{1}_m &= \nu
\end{aligned}
\right\}.
\]

\begin{theorem}[OT dominates mean-pooling CLIP score]
\label{thm:ot_ge_mean}
Assume CLIP aggregates segment/token features by mean pooling:
\[
F_M=\frac{1}{m}\sum_{i=1}^m f_{m_i},
\qquad
F_T=\frac{1}{n}\sum_{j=1}^n f_{t_j},
\]
and uses the global score
\[
R_{\mathrm{CLIP}}
\;=\;
\langle F_M, F_T\rangle .
\]
Define the OT reward (fine-grained score) as
\[
\begin{aligned}
R_{\mathrm{OT}}
&=
\max_{\pi\in\Pi(\mu,\nu)}
\;\langle \pi, C\rangle \\
&=
\max_{\pi\in\Pi(\mu,\nu)}
\sum_{i=1}^m
\sum_{j=1}^n
\pi_{ij} C_{ij}.
\end{aligned}
\]
Then for any fixed features,
\[
R_{\mathrm{OT}}
\;\ge\;
R_{\mathrm{CLIP}}.
\]
Equivalently, defining costs by
$d_{ij}=-C_{ij}$ and
\[
\mathrm{Cost}_{\mathrm{OT}}
=
\min_{\pi\in\Pi(\mu,\nu)}
\langle \pi, d\rangle ,
\]
we have
\[
\mathrm{Cost}_{\mathrm{OT}}
\;\le\;
\langle U, d\rangle ,
\qquad
U_{ij}=\frac{1}{mn},
\]
i.e., OT achieves a cost no larger than the
mean-pooling induced plan.
\end{theorem}

\begin{proof}
By bilinearity,
\[
\begin{aligned}
R_{\mathrm{CLIP}}
&=
\left\langle
\frac{1}{m}\sum_{i=1}^m f_{m_i},
\frac{1}{n}\sum_{j=1}^n f_{t_j}
\right\rangle \\
&=
\frac{1}{mn}
\sum_{i=1}^m
\sum_{j=1}^n
\langle f_{m_i}, f_{t_j}\rangle \\
&=
\sum_{i=1}^m
\sum_{j=1}^n
U_{ij} C_{ij}
=
\langle U, C\rangle ,
\end{aligned}
\]
where $U\in\mathbb{R}^{m\times n}$ is defined by
$U_{ij}=\frac{1}{mn}$.

We verify feasibility.
For each $i$,
\[
\sum_{j=1}^n U_{ij}
=
n\cdot\frac{1}{mn}
=
\frac{1}{m}
=
\mu_i,
\]
and for each $j$,
\[
\sum_{i=1}^m U_{ij}
=
m\cdot\frac{1}{mn}
=
\frac{1}{n}
=
\nu_j.
\]
Hence $U\in\Pi(\mu,\nu)$.

Since
$R_{\mathrm{OT}}
=
\max_{\pi\in\Pi(\mu,\nu)}
\langle \pi, C\rangle$
and $U$ is feasible,
\[
R_{\mathrm{OT}}
\;\ge\;
\langle U, C\rangle
=
R_{\mathrm{CLIP}}.
\]
The cost-form statement follows by
$d=-C$ and
\[
\min_{\pi}\langle \pi, d\rangle
=
-\max_{\pi}\langle \pi, C\rangle .
\]
\end{proof}

\paragraph{Frozen-encoder retrieval setting.}
Let the training set be
$\mathcal{D}=\{(M_i,T_i)\}_{i=1}^N$.
Fix (freeze) the encoders, so for any pair $(M,T)$
the feature sets and
$C(M,T)$ are deterministic.

Given a score function $S(M,T)$,
define top-1 retrieval by
\[
\begin{aligned}
\widehat{T}(M_q)
&=
\arg\max_{T_k\in\{T_1,\dots,T_N\}}
S(M_q,T_k), \\
\mathrm{Acc}(S;\mathcal{D})
&=
\frac{1}{N}
\sum_{q=1}^N
\mathbf{1}\{\widehat{T}(M_q)=T_q\}.
\end{aligned}
\]

Let the CLIP (mean-pooling) score be
\[
S_{\mathrm{mean}}(M,T)
=
\langle U, C(M,T)\rangle .
\]
Consider the fine-grained (transport) family
\[
\mathcal{H}_{\mathrm{OT}}
=
\left\{
S_{\pi}(M,T)
=
\langle \pi, C(M,T)\rangle
\;\middle|\;
\pi\in\Pi(\mu,\nu)
\right\}.
\]
Note that
$\mathcal{H}_{\mathrm{OT}}$
includes the mean-pooling scorer since
$U\in\Pi(\mu,\nu)$.

\begin{theorem}[Fine-grained OT family has no worse achievable training retrieval accuracy]
\label{thm:acc_superset}
With frozen encoders and a fixed dataset
$\mathcal{D}$,
\[
\max_{S\in\mathcal{H}_{\mathrm{OT}}}
\mathrm{Acc}(S;\mathcal{D})
\;\ge\;
\mathrm{Acc}(S_{\mathrm{mean}};\mathcal{D}).
\]
Moreover, if there exists a query $M_q$
and a negative $T_k\neq T_q$ such that
\[
S_{\mathrm{mean}}(M_q,T_k)
=
S_{\mathrm{mean}}(M_q,T_q),
\]
but there exists
$\pi^\star\in\Pi(\mu,\nu)$
with
\[
S_{\pi^\star}(M_q,T_q)
>
S_{\pi^\star}(M_q,T_k),
\]
then the inequality is strict:
\[
\max_{S\in\mathcal{H}_{\mathrm{OT}}}
\mathrm{Acc}(S;\mathcal{D})
\;>\;
\mathrm{Acc}(S_{\mathrm{mean}};\mathcal{D}).
\]
\end{theorem}

\begin{proof}
Since $U\in\Pi(\mu,\nu)$,
we have
$S_{\mathrm{mean}} = S_U \in \mathcal{H}_{\mathrm{OT}}$.
Therefore,
\[
\max_{S\in\mathcal{H}_{\mathrm{OT}}}
\mathrm{Acc}(S;\mathcal{D})
\;\ge\;
\mathrm{Acc}(S_U;\mathcal{D})
=
\mathrm{Acc}(S_{\mathrm{mean}};\mathcal{D}).
\]

For strictness,
assume there exists $(q,k)$ such that
\[
S_{\mathrm{mean}}(M_q,T_k)
=
S_{\mathrm{mean}}(M_q,T_q),
\]
but some
$\pi^\star$
satisfies
\[
S_{\pi^\star}(M_q,T_q)
>
S_{\pi^\star}(M_q,T_k).
\]
Under $S_{\pi^\star}$,
the positive ranks above the negative for query $M_q$,
resolving at least one ambiguity
that mean pooling cannot.
Hence there exists a scorer in
$\mathcal{H}_{\mathrm{OT}}$
with strictly higher top-1 accuracy on
$\mathcal{D}$.
\end{proof}
